\documentclass[10pt,a4paper]{amsart}
\usepackage[margin=19mm]{geometry}
\usepackage{amsmath,amssymb,mathtools}
\usepackage[scr=rsfs,cal=euler]{mathalfa}
\usepackage{xfrac}
\usepackage[unicode,hidelinks,bookmarks=false]{hyperref}
\newcommand{\ie}{i.e.\ }
\newcommand{\cf}{cf.\ }
\newcommand{\resp}{respectively\ }
\newcommand{\Subsection}{\subsection}
\providecommand{\nobreakdash}{\nobreak\hbox{-}}
\setlength{\emergencystretch}{2em}
\multlinegap0pt
\usepackage{bm}
\usepackage{bbm}
\usepackage{dsfont}
\usepackage{xspace}
\usepackage{cancel}
\usepackage{mathtools}
\usepackage{cleveref}
\usepackage{amsthm}
\makeatletter
\@ifundefined{theorem}{\newtheorem{theorem}{Theorem}}{}
\makeatother
\title[Enumerating pattern-avoiding permutations by leading terms]{Enumerating pattern-avoiding permutations by leading terms}
\author{\"Omer E\u{g}ecio\u{g}lu, Collier Gaiser, Mei Yin}
\date{Source: arXiv:2309.15964v4 (CC BY 4.0)}
\begin{document}
\maketitle

\noindent\emph{Source and licence.} Enumerating pattern-avoiding permutations by leading terms. Version arXiv:2309.15964v4 (CC BY 4.0), distributed under the Creative Commons Attribution 4.0 International licence, as recorded in the arXiv licence metadata for this exact version and shown on the abstract page https://arxiv.org/abs/2309.15964v4. Full rights notice: licenses/arxiv-2309.15964-CC-BY-4.0.txt.\par\smallskip

\noindent\textit{Editorial scope.} Counted unit: the Theorem whose source label is \texttt{None} in the article (source0.tex lines 814--828, source label \texttt{None}), delivered together with the article's own complete original proof of it (source0.tex lines 829--865, one \texttt{proof} environment of 37 lines). No mathematics is written, repaired or re-derived: statement and proof are the article's own text line for line. Required context retained uncounted: Theorem:34123421 (lines 738--745). Every definition, display and label of those lines is retained unchanged. Omitted from this package and not redistributed: 1--737, 746--813, 866--1169. Nothing inside a retained excerpt is omitted or altered: every retained body is delivered exactly as the article's lines are written, so no omission receipt of any kind accompanies this package. Citations: the retained excerpts carry no citation command at all, so no bibliography entry of the article is transcribed and no reference is redistributed. Reference adaptation: none was needed and none was made; every reference inside the delivered unit resolves inside it, so no unresolved reference or citation is produced. Printed slips kept literal: no printed word, symbol, hypothesis or number of the counted statement or of its proof was altered. Layout and class: the article's own class is \texttt{article} with packages including amssymb, float, amsthm, hyperref, xcolor, amsmath, tipa, witharrows, cleveref, geometry; the wrapper is the lane's pinned \texttt{amsart} editorial wrapper at 10pt on A4, which restates the theorem-like environments the retained text needs and adds the article's own macro definitions that the retained text needs (none), with extra wrapper packages (bm, bbm, dsfont, xspace, cancel, mathtools, cleveref). The delivered numbering is the unit's own. Assets: no retained span contains a figure environment, a graphics-inclusion command, a PDF-page inclusion, a table or a TikZ picture; no image file is copied into this package, the package declares no asset dependency and no image file is redistributed with it. Licence: version arXiv:2309.15964v4 of this article is distributed under the Creative Commons Attribution 4.0 International licence, as recorded in the arXiv licence metadata for this exact version and shown on the abstract page https://arxiv.org/abs/2309.15964v4. A full rights notice is delivered at licenses/arxiv-2309.15964-CC-BY-4.0.txt. Characters: every retained excerpt is pure ASCII, so the delivered engine, XeLaTeX with its default Latin Modern text font, reports no missing character.
\begin{theorem}\label{Theorem:34123421}

For all $n\geq1$,
\[
|S_n(3412,3421)|=\mathbb{S}_{n-1},
\]
where $\mathbb{S}_{n-1}$ is the $(n-1)$st large (big) Schr\"{o}der number.
\end{theorem}

\begin{theorem}
    Let
    \[
    U=\{c_i:i\in[t]~\text{and there exist}~j,k\in[t]~\text{such that}~i<j<k~\text{and}~c_ic_jc_k~\text{is a 231 pattern}\}
    \]
    and
    \[
    V=\{c_i:i\in[t]~\text{and there exists}~j\in[t]~\text{such that}~i<j~\text{and}~c_i<c_j\}.
    \]
    If $U\neq\emptyset$ and $\left|[\max U]\backslash\{c_1,c_2,\ldots,c_t\}\right|\geq1$ or $V\neq\emptyset$ and $\left|[\max V]\backslash\{c_1,c_2,\ldots,c_t\}\right|\geq2$, then $|S_{n,(c_1,c_2,\ldots,c_t)}(3412,3421)|=0$; otherwise,
\[
|S_{n,(c_1,c_2,\ldots,c_t)}(3412,3421)|=\mathbb{S}_{c_{(j)}-c_{(j-1)}-2}\prod_{i=j+1}^{t+1}\mathbb{S}_{c_{(i)}-c_{(i-1)}-1},
\]
where $c_{(0)}=0$, $c_{(t+1)}=n+1$, $c_{(1)}<c_{(2)}<\cdots<c_{(t)}$ are the order statistics of $\{c_1,c_2,\ldots,c_t\}$, and $j=\min\{i\in[t+1]:c_{(i)}-c_{(i-1)}>1\}$.
\end{theorem}

\begin{proof}
    For all $k\in[t+1]$, let $A_k=\{c_{(k-1)}+1,c_{(k-1)}+2,\ldots,c_{(k)}-1\}$. We note that it is possible that $A_k=\emptyset$ for some $k$. Also notice that
    \[
[n]\backslash\{c_1,c_2,\ldots,c_t\}=\bigcup_{k=1}^{t+1}A_k.
\]


We first suppose that $U\neq\emptyset$ and $\left|[\max U]\backslash\{c_1,c_2,\ldots,c_t\}\right|\geq1$. Let $x,y,z\in[t]$ such that $c_x=\max U$, $x<y<z$, and $c_xc_yc_z$ is a $231$ pattern. Since $\left|[c_x]\backslash\{c_1,c_2,\ldots,c_t\}\right|\geq1$, there exists $\alpha\in [n]\backslash\{c_1,c_2,\ldots,c_t\}$ such that $\alpha<c_x$. If $\alpha<c_z$, then $c_xc_yc_z\alpha$ is a $3421$ pattern; and if $\alpha>c_z$, then $c_xc_yc_z\alpha$ is a $3412$ pattern. Hence $|S_{n,(c_1,c_2,\ldots,c_t)}(3412,3421)|=0$.

Next, we suppose that $V\neq\emptyset$ and $\left|[\max V]\backslash\{c_1,c_2,\ldots,c_t\}\right|\geq2$. Let $x,y\in[t]$ such that $c_x=\max V$, $x<y$, and $c_x<c_y$. Since $\left|[c_x]\backslash\{c_1,c_2,\ldots,c_t\}\right|\geq2$, there exist $\alpha,\beta\in[n]\backslash\{c_1,c_2,\ldots,c_t\}$ such that $\alpha<\beta<c_x<c_y$. If $\alpha\in\mathcal{A}_\tau(\beta)$, then $c_xc_y\alpha\beta$ is a $3412$ pattern; and if $\alpha\in\mathcal{D}_\tau(\beta)$, then $c_xc_y\beta\alpha$ is a $3421$ pattern. Hence $|S_{n,(c_1,c_2,\ldots,c_t)}(3412,3421)|=0$.

Now suppose otherwise. Let $j=\min\{i\in[t+1]:c_{(i)}-c_{(i-1)}>1\}$. In other words, $j\in[t+1]$ is the smallest index such that $A_j\neq\emptyset$. For all $\tau\in S_{n,(c_1,c_2,\ldots,c_t)}$, let $\tau_j$ be the subpermutation of $\tau$ on $A_j$, and for all $i\in\{j+1,j+2,\ldots,t+1\}$, let $x_i$ be the last term of the subpmermutation of $\tau$ on $A_j\cup A_{j+1}\cup\cdots \cup A_{i-1}$ and $\tau_i$ be the subpermutation of $\tau$ on $A_i\cup\{x_i\}$.
    
    Let $\mathcal{R}'$ be the subset of $S_{n,(c_1,c_2,\ldots,c_t)}$ such that every $\tau\in\mathcal{R}'$ satisfies the following:
    \begin{itemize}
        \item[(i)] for all $i\in\{j+1,j+2,\ldots,t+1\}$, $y\in A_i$, and $z\in (A_j\cup A_{j+1}\cup\cdots \cup A_{i-1})\backslash\{x_i\}$, we have $z\in\mathcal{A}_\tau(y)$;
        \item[(ii)] and for all $i\in\{j,j+1,\ldots,t+1\}$, $\tau_i$ avoids both $3412$ and $3421$.
    \end{itemize}
    
    Let $\tau\in\mathcal{R}'$. By \cref{Theorem:34123421}, since $|A_j|=c_{(j)}-c_{(j-1)}-1$, there are $\mathbb{S}_{c_{(j)}-c_{(j-1)}-2}$ possibilities for $\tau_j$. Now let $i\in\{j+1,j+2,\ldots,t+1\}$. Inductively, we count the possibilities for $\tau_i$ when $\tau_j,\tau_{j+1},\ldots,\tau_{i-1}$ are fixed. In this case, since $|A_i\cup\{x_i\}|=c_{(i)}-c_{(i-1)}$, by \cref{Theorem:34123421}, there are $\mathbb{S}_{c_{(i)}-c_{(i-1)}-1}$ possibilities for $\tau_i$. Hence, we have $|\mathcal{R}'|=\mathbb{S}_{c_{(j)}-c_{(j-1)}-2}\prod_{i=j+1}^{t+1}\mathbb{S}_{c_{(i)}-c_{(i-1)}-1}$.

    It remains to show that $\mathcal{R}'=S_{n,(c_1,c_2,\ldots,c_t)}(3412,3421)$. Let $\tau\in S_{n,(c_1,c_2,\ldots,c_t)}(3412,3421)$. Then Property (ii) is obviously satisfied. Now we show that $\tau$ satisfies Property (i). Suppose, by way of contradiction, that $\tau$ does not satisfy Property (i). Then there exist $i\in\{j+1,j+2,\ldots,t+1\}$, $y\in A_i$, and $z\in (A_j\cup A_{j+1}\cup\cdots \cup A_{i-1})\backslash\{x_i\}$ such that $z\in D_\tau(y)$. Since $x_i$ is the last term of the subpermutation of $\tau$ on $A_j\cup A_{j+1}\cup\cdots\cup A_{i-1}$, $c_{(i-1)}yzx_i$ is a subpermutation of $\tau$. Since $y\in A_i$, we have $y>c_{(i-1)}$. Now since $z,x_i\in A_j\cup A_{j+1}\cup\cdots \cup A_{i-1}$, $c_{(i-1)}yzx_i$ is either a $3412$ pattern or a $3421$ pattern. This is a contradiction. Hence $\tau$ satisfies Property (i). It follows that $S_{n,(c_1,c_2,\ldots,c_t)}(3412,3421)\subseteq\mathcal{R}'$. 

    We still need to prove that $\mathcal{R}'\subseteq S_{n,(c_1,c_2,\ldots,c_t)}(3412,3421)$. Let $\tau\in \mathcal{R}'$. Suppose, by way of contradiction, that $abcd$ is a subpermutation of $\tau$ which is either a $3412$ pattern or a $3421$ pattern. Then we have $c<d<a<b$ or $d<c<a<b$. We have five cases:

    Case 1: $a,b,c,d\notin\{c_1,c_2,\ldots,c_t\}$. Then $a\in A_k$ for some $k\in\{j,j+1,\ldots,t+1\}$. Since $c,d<a$, we have $c\in A_{i_1}$ and $d\in A_{i_2}$ for some $i_1,i_2\leq k$. Since $b>a$, if $b\notin A_k$, then $b\in A_{i_3}$ with $i_3>k$. Then Property (i) is violated because $c,d\in A_j\cup A_{j+1}\cup\cdots\cup A_{k}$. So we must have $b\in A_k$. By Property (i), at most one of $c$ and $d$ is in $A_j\cup A_{j+1}\cup\cdots\cup A_{k-1}$. In addition, if $c$ or $d$ is in $A_j\cup A_{j+1}\cup\cdots\cup A_{k-1}$, then it must be the last term of $\tau_{k-1}$. So $abcd$ is a subpermutation on $A_k\cup\{x_k\}$. This contradicts Property (ii).
    
    Case 2: $a\in\{c_1,c_2,\ldots,c_t\}$ but $b,c,d\notin\{c_1,c_2,\ldots,c_t\}$. Then $a=c_k$ for some $k\in[t]$. Since $b>a$, we have $b\in A_{i}$ for some $i>k$. Since $c,d<a$, we have $c,d\in A_j\cup A_{j+1}\cup\cdots \cup A_{k}$. Since $c,d\in D_\tau(b)$, this violates Property (i). Hence we have a contradiction.

    Case 3: $a,b\in\{c_1,c_2,\ldots,c_t\}$ but $c,d\notin\{c_1,c_2,\ldots,c_t\}$. Since $a<b$ and $c,d<a$, $V\neq\emptyset$ and $\left|[\max V]\backslash\{c_1,c_2,\ldots,c_t\}\right|\geq2$ which is a contradiction.

    Case 4: $a,b,c\in\{c_1,c_2,\ldots,c_t\}$ but $d\notin\{c_1,c_2,\ldots,c_t\}$. Then $abc$ is a $231$ pattern. Since $d<a$, $U\neq\emptyset$ and $\left|[\max U]\backslash\{c_1,c_2,\ldots,c_t\}\right|\geq1$ which is a contradiction.

    Case 5: $a,b,c,d\in\{c_1,c_2,\ldots,c_t\}$. This contradicts our convention that $(c_1,c_2,\ldots,c_t)$ avoids both $3412$ and $3421$.

    Hence, $\tau$ avoids both $3412$ and $3421$. It follows that $\mathcal{R}'\subseteq S_{n,(c_1,c_2,\ldots,c_t)}(3412,3421)$. This completes our proof.
\end{proof}

\end{document}
